\section{Structures, what proofs use, and changes of description}\label{appE:sec:structures}
Almost every object in this book is a set with something extra attached, which Section~\ref{ch06:sec:structure} calls a \emph{structure}. It may be a \emph{relation} (Section~\ref{ch06:sec:differences}), which links or compares pairs of objects, as congruence modulo $m$ and adjacency in a graph do, or an \emph{operation}, such as addition, which combines two elements of a set into an element of the same set. Functions, also called maps (Section~\ref{ch03:sec:functions}), let us carry a question from one description to another.

The extra structure decides which questions can even be asked.
\begin{itemize}
\item Addition and scaling let us form linear combinations of vectors (Chapter~\ref{ch:linear}).
\item Order lets us pick out a smallest or a largest object. Every nonempty set of natural numbers has a smallest member (the well-ordering property of Sections~\ref{ch04:sec:contrapositive} and~\ref{ch05:sec:not-circular}), and every nonempty finite one has a largest; this is how Chapter~\ref{ch:graphs} chooses a shortest walk between two vertices joined by some walk, and a longest directed path in a nonempty finite directed graph.
\item A distance lets us say that one point approximates another, and so speak of limits (Chapter~\ref{ch:limits}). Without one, ``$x_n$ gets close to $x_*$'' has no meaning.
\item A probability on a finite set of outcomes lets us take weighted averages (Chapter~\ref{ch:probability}).
\end{itemize}

Habit~3 (Section~\ref{ch03:sec:images}) asks which information a proof actually uses. For the main theorems of the book the answers are short.
\begin{itemize}
\item The proofs of Hall's theorem (Chapter~\ref{ch:hall}) use only which students accept which tasks and the sizes of finite sets, never what the students and tasks are. So the theorem applies to every finite bipartite graph, whatever its vertices stand for.
\item The proof of the contraction theorem (Theorem~\ref{thm:banach}) uses, besides the facts that the space is nonempty and that $T$ maps it into itself, only the contraction inequality, completeness, and the three requirements of a metric, above all the triangle inequality. It never adds points or multiplies them by numbers, so it works in metric spaces that have no addition at all.
\item The proof of the Cauchy--Schwarz inequality (Theorem~\ref{thm:cs}) uses only three properties of the dot product: linearity in each argument, symmetry, and positivity. So the same proof works for any other way of multiplying two vectors to get a number, provided it has these three properties and the length of $u$ is defined as the square root of the product of $u$ with itself. For example, on $\RR^2$ the product $2u_1v_1+u_2v_2$ is symmetric and linear in each argument, and positive, since $2u_1^2+u_2^2$ is never negative and is zero only when $u=0$. So $|2u_1v_1+u_2v_2|\le\sqrt{2u_1^2+u_2^2}\,\sqrt{2v_1^2+v_2^2}$ for all $u,v\in\RR^2$, with no new work.
\item Chapter~\ref{ch:signs} turns the condition that a sign sequence has zero periodic correlation at every nonzero shift into the condition that different rows of its circulant matrix are orthogonal (Proposition~\ref{ch14:prop:rows}). Adding up all the correlations shows that the order of the matrix, its number of rows, is a perfect square (Section~\ref{ch14:sec:square}), an argument that relies on the circulant pattern. Comparing three rows shows that an order greater than $2$ is divisible by four (Lemma~\ref{thm:hadamardorder}). That argument uses only that the entries are $\pm1$ and different rows are orthogonal, which is exactly what makes a square matrix a Hadamard matrix (Section~\ref{ch14:sec:circulant}), so it holds for every Hadamard matrix, circulant or not.
\end{itemize}
Noticing what a proof does \emph{not} use tells you where its conclusion might reach further. This is a guide for guessing, not a proof: every proposed extension must still be checked against the argument, line by line.

Many steps in the book change the description of a problem without changing the problem, and Habit~6 (Section~\ref{ch06:sec:order}) asks what such a change preserves; a new description helps only when we know exactly what it keeps and what it throws away. Here are four from the book.
\begin{itemize}
\item \emph{An inverse loses nothing.} If $f$ has an inverse $h$, every input is recovered from its output as $x=h(f(x))$. In particular, $f(x)=f(y)$ gives $x=h(f(x))=h(f(y))=y$, so different inputs never share an output, and a question about inputs can be answered by looking at outputs (Section~\ref{ch03:sec:injsurj}).
\item \emph{A quotient forgets on purpose.} Passing from an integer $a$ to its residue class $[a]_m$ forgets everything except the remainder on division by $m$. This is safe only for rules that give the same answer for every representative of a class, which is why well-definedness has to be checked (Section~\ref{ch06:sec:well-defined}).
\item \emph{A matrix turns a linear map into a table of numbers.} Column $j$ of the matrix records where the map sends the standard basis vector $e_j$ (Section~\ref{ch10:sec:linearmaps}), and composing linear maps corresponds to multiplying their matrices (Section~\ref{ch10:sec:product}).
\item \emph{An indicator turns an event into a quantity.} An indicator is $1$ when its event occurs and $0$ when it does not, so counting how many events occur becomes adding indicators, and linearity of expectation applies (Section~\ref{ch13:sec:indicators}).
\end{itemize}
Each of them gives a language in which the step we need is easier to state and to justify.
